# SimPath: Technical Specification v2.0
# (Reviewer-Hardened Edition)

> **Full title**: SimPath: Simulation-in-the-Loop Learning Path Recommendation
> via Robust Student Persona Testing
> **Target venues**: AIED 2026 Main Track / RecSys 2026
> **Version**: 2.0 — March 2026
> **Changelog from v1**: Reviewer defense strategies integrated throughout.
> Model choices updated. Circular evaluation fix. Persona validation added.
> Retrieve-then-order path generation. Terminology adjusted.

---

## 0. Notation Table

| Symbol | Type | Description |
|--------|------|-------------|
| $\mathcal{S} = \{s_1, \dots, s_N\}$ | set | All students in dataset |
| $\mathcal{Q} = \{q_1, \dots, q_M\}$ | set | All questions/exercises |
| $\mathcal{K} = \{k_1, \dots, k_C\}$ | set | Knowledge components (KC) |
| $\mathbf{h}_n$ | sequence | Interaction history of student $n$: $[(q_1, r_1, t_1), \dots]$; $r_i \in \{0,1\}$, $t_i$ = timestamp |
| $\mathbf{m}_n \in [0,1]^C$ | vector | Knowledge state of student $n$; $m_n^c$ = mastery of KC $c$ |
| $P = [q_{i_1}, \dots, q_{i_L}]$ | sequence | Learning path of length $L$ |
| $\mathcal{P} = \{P_1, \dots, P_K\}$ | set | $K$ candidate paths |
| $\Theta = \{\theta_{\text{grit}}, \theta_{\text{frag}}, \theta_{\text{real}}\}$ | set | Three persona simulators |
| $\phi_j(P_i)$ | scalar | Composite score of path $P_i$ under persona $\theta_j$ |
| $R(P_i, \theta_j)$ | scalar | Regret of path $P_i$ under persona $\theta_j$ |

**Terminology note**: We use "persona simulator" (not "agent") for the
student models. The term "agent" is reserved for the LLM components that
perform generation and explanation. This distinction avoids inflated claims
about agent autonomy. [Defense: Reviewer #8]

---

## 1. Key Design Decisions (Reviewer-Driven)

This section documents architectural choices made specifically to preempt
known reviewer concerns. Each links to the attack point it addresses.

### 1.1 Separate KT Models for Recommendation vs Evaluation [Defense: #1]

```
KT_recommend = AKT   (used in path generation + persona simulation)
KT_evaluate  = SAINT (used only for offline evaluation metrics)

Rationale: If the same KT model generates recommendations AND evaluates them,
any model bias appears as inflated performance. Using two structurally different
KT models (attention-based AKT vs transformer-based SAINT) breaks this
circularity. Additionally, we report correlation with held-out ground truth
(see Section 7.2).
```

### 1.2 Retrieve-Then-Order Path Generation [Defense: #5]

```
Instead of letting the LLM freely generate question_ids (hallucination risk),
we split path generation into:
  Step A: RETRIEVAL — deterministic filtering of question pool based on
          KC mastery, difficulty range, staleness. Pure code, no LLM.
  Step B: ORDERING — LLM receives the filtered pool and decides the
          optimal sequence. LLM can only SELECT from existing IDs.

This makes hallucination structurally impossible. The LLM cannot invent
question_ids because it only sees real ones.
```

### 1.3 Persona Validation via Data Clustering [Defense: #2]

```
Before running experiments, validate persona definitions:
  1. Extract behavioral features for ALL students in training set:
     - avg_session_length, dropout_after_k_wrong (k=1..5),
       avg_elapsed_time, skip_rate, overall_accuracy
  2. Run k-means clustering (k=3) on standardized features
  3. Compare cluster centroids to θ_grit, θ_frag, θ_real parameters
  4. Report cosine similarity between centroids and persona definitions
  5. Visualize with t-SNE / UMAP — include as Figure in paper

This provides empirical evidence that our 3 personas correspond to
naturally occurring student subpopulations, not arbitrary constructs.
```

### 1.4 Persona = Session State, Not Student Type [Defense: #4]

```
CRITICAL FRAMING (use in Introduction + Related Work):

Personas do NOT represent fixed student types.
The same student can be "Gritty" on Monday morning (well-rested,
motivated) and "Fragile" on Friday night (tired, distracted).

Therefore, personas model the UNCERTAINTY OF SESSION-LEVEL BEHAVIOR,
not the uncertainty of student identity. This justifies minimax regret:
we don't know which "version" of the student will show up today.

Supporting citation: Pekrun (2006) — emotions and motivation fluctuate
within individuals across learning sessions.
```

### 1.5 Model Choices [Defense: #6]

```
PRIMARY MODELS (main experiments):
  - GPT-5.4         (for path ordering + explanation generation)
  - Claude Sonnet 4.6 (for path ordering + explanation generation)

Both are used for ALL main experiments. Every result table reports both.
This demonstrates model-agnostic performance across two different LLM
families (OpenAI vs Anthropic).

ABLATION (secondary):
  - Open-source models (Llama-3.1-8B, Qwen2.5-7B) as future extension
  - Mentioned in Discussion as planned reproducibility experiment

LOGGING:
  - Record exact model version strings (e.g., "gpt-5.4-2026-03-05",
    "claude-sonnet-4-6")
  - Log all prompts and responses for reproducibility
  - Temperature, seed, max_tokens fixed and reported
```

### 1.6 Agent4Edu Differentiation [Defense: #7]

```
Use this framing consistently throughout the paper:

  Agent4Edu (AAAI 2025) = "simulation FOR evaluation"
    → Generates synthetic student data to evaluate existing systems
    → Simulation is the END PRODUCT

  SimPath (Ours) = "simulation FOR recommendation"
    → Simulates student personas to STRESS-TEST candidate paths
    → Simulation is an INTERMEDIATE STEP in the recommendation pipeline
    → The end product is a ROBUST RECOMMENDATION + EXPLANATION

Additionally, SimPath's persona simulators are lightweight (KT + params),
not full LLM agents like Agent4Edu. This is a deliberate design choice
for cost efficiency (see Section 9).

Include comparison table in Related Work:

| Aspect         | Agent4Edu          | SimPath (Ours)        |
|----------------|--------------------|-----------------------|
| Purpose        | Data generation    | Path recommendation   |
| Simulation     | End product        | Intermediate step     |
| Architecture   | LLM agent (heavy)  | KT + params (light)   |
| Output         | Synthetic data     | Recommended path      |
| Selection      | N/A                | Minimax regret        |
| Explainability | N/A                | Simulation-grounded   |
```

---

## 2. Data Preprocessing Pipeline

### 2.1 Dataset Selection

**Primary: EdNet KT3**
- Source: https://github.com/riiid/ednet
- Level: KT3 (question-response + explanation + elapsed time)
- Size: ~297K students, 13,169 questions, 293 skill tags
- Domain: TOEIC English (listening + reading)

**Secondary: ASSISTments 2015**
- Source: https://sites.google.com/site/assistmaborern/
- Size: ~19,917 students, 100 skills
- Domain: Middle school mathematics

**Domain diversity justification [Defense: #9]**: Using two completely
different domains (English testing vs math) demonstrates that SimPath is
domain-agnostic. If results are consistent across both, this is STRONGER
evidence than using two math datasets.

### 2.2 Preprocessing Steps

```
INPUT:  raw interaction logs
OUTPUT: cleaned sequences + behavioral features + question metadata

Step 1: Filter students with < 20 or > 5000 interactions
Step 2: Extract per-interaction features:
        - question_id, skill_tags (KC ids), correctness (0/1)
        - elapsed_time (seconds), timestamp (unix)
Step 3: Sort by timestamp within each student
Step 4: Compute per-question difficulty:
        difficulty(q) = 1 - (# correct / # total responses)
Step 5: Split each student's sequence:
        - First 70%  → KT training set
        - Next 10%   → KT validation set
        - Next 10%   → SimPath recommendation input (recent history)
        - Last 10%   → Held-out ground truth (NEVER used by SimPath)
                        [Defense: #1 — separate evaluation data]
Step 6: Build skill-question mapping
Step 7: Extract per-student behavioral features (for Realistic persona +
        for persona validation clustering):
        - avg_session_length (gap > 30min = new session)
        - dropout_pattern: P(leave | consec_wrong >= k) for k=1..5
        - avg_elapsed_time per question
        - skip_rate: fraction answered in < 5 seconds
        - accuracy_by_difficulty_bin: [easy, medium, hard]
```

### 2.3 Persona Validation Clustering [Defense: #2]

```
PROCEDURE validate_personas(all_students):
  # Extract behavioral feature vectors
  features = []
  FOR student_n in all_students:
    f = [
      avg_session_length(n),
      dropout_rate_at_k2(n),     # P(leave | 2 consec wrong)
      avg_elapsed_time(n),
      skip_rate(n),
      overall_accuracy(n)
    ]
    features.append(standardize(f))

  # Cluster
  kmeans = KMeans(n_clusters=3, random_state=42, n_init=10)
  labels = kmeans.fit_predict(features)
  centroids = kmeans.cluster_centers_

  # Compare to persona definitions
  θ_grit_vec = standardize([high_session, low_dropout, high_time, low_skip, high_acc])
  θ_frag_vec = standardize([low_session, high_dropout, low_time, high_skip, low_acc])
  θ_real_vec = centroids[largest_cluster]  # most common type

  FOR i, centroid in enumerate(centroids):
    print(f"Cluster {i}: cosine_sim to Gritty={cosine(centroid, θ_grit_vec):.3f}, "
          f"Fragile={cosine(centroid, θ_frag_vec):.3f}")

  # Visualize with UMAP
  reducer = UMAP(n_components=2, random_state=42)
  embedding = reducer.fit_transform(features)
  # Plot with cluster colors → Figure 2 in paper

  RETURN centroids, labels, cosine_similarities
```

**Expected output**: One cluster aligns with Gritty (high persistence,
long sessions), one with Fragile (high dropout, short sessions), one in
between (Realistic). Report cosine similarities in paper.

---

## 3. Knowledge Tracing Models (Phase 1)

### 3.1 Two-Model Setup [Defense: #1]

```
MODEL A: AKT (Attentive Knowledge Tracing)
  - Used for: path generation, persona simulation
  - Why: strong per-KC mastery estimation, attention mechanism

MODEL B: SAINT (Separated Attentive INternet-based Tracing)
  - Used for: offline evaluation ONLY
  - Why: transformer encoder-decoder, structurally different from AKT
  - This separation breaks circular evaluation

Both models trained on the SAME training data (Step 5 first 70%).
Both validated on the SAME validation data (Step 5 next 10%).
```

### 3.2 Training Hyperparameters

```
AKT:
  embedding_dim=256, num_heads=8, num_blocks=4, dropout=0.1
  lr=1e-3, batch_size=64, max_seq_len=200, epochs=100
  early_stopping: patience=10 on validation AUC

SAINT:
  d_model=256, num_heads=8, enc_layers=4, dec_layers=4, dropout=0.1
  lr=1e-3, batch_size=64, max_seq_len=200, epochs=100
  early_stopping: patience=10 on validation AUC

Loss: Binary cross-entropy
  L = -(1/T) * Σ [r_t * log(p_t) + (1-r_t) * log(1-p_t)]
```

### 3.3 Knowledge State Extraction

```
FUNCTION extract_knowledge_state(kt_model, history_n):
  """
  INPUT:  trained KT model, student interaction history
  OUTPUT: mastery vector m_n ∈ [0,1]^C
  """
  Feed history_n through kt_model

  FOR each KC c:
    last_interaction_with_c = find_last(history_n, kc=c)
    IF last_interaction_with_c exists:
      m_n[c] = kt_model.predict_mastery(history_n, c, at=last_interaction)
    ELSE:
      # Never encountered: use population prior
      m_n[c] = population_mean_accuracy(c)

  RETURN m_n
```

### 3.4 KT Simulation Function

```
FUNCTION simulate_kt(kt_model, m_n, path P, persona_params θ):
  """
  Simulates a student walking through path P under persona θ.
  Returns trajectory, final mastery, and dropout status.

  NOTE: This uses kt_model = AKT (recommendation model).
  Evaluation uses a SEPARATE model (SAINT). [Defense: #1]
  """
  m_current = m_n.copy()
  trajectory = []
  dropout = False
  consecutive_wrong = 0

  FOR step t = 1 to |P|:
    q_t = P[t]
    kc_t = skill_tags(q_t)
    diff_t = difficulty(q_t)

    # --- Skip check ---
    IF θ.skip_threshold is not None AND diff_t > θ.skip_threshold:
      IF Bernoulli(0.5):
        trajectory.append({q: q_t, result: "skip", m: m_current.copy()})
        consecutive_wrong += 1  # skip counts as failure for dropout
        GOTO dropout_check

    # --- Predict correctness ---
    p_correct = kt_model.predict(m_current, q_t)
    p_adjusted = clip(p_correct + θ.p_correct_boost, 0.01, 0.99)
    r_t = Bernoulli(p_adjusted)

    # --- Update mastery ---
    FOR each kc in kc_t:
      IF r_t == 1:
        m_current[kc] += θ.learn_rate * (1 - m_current[kc])
        consecutive_wrong = 0
      ELSE:
        m_current[kc] -= θ.forget_penalty * m_current[kc] * 0.1
        consecutive_wrong += 1

    trajectory.append({q: q_t, result: r_t, m: m_current.copy()})

    # --- Dropout check ---
    dropout_check:
    p_dropout = min(1.0, θ.dropout_base(consecutive_wrong) *
                         (1 + θ.difficulty_sensitivity * diff_t))
    IF Bernoulli(p_dropout):
      dropout = True
      BREAK

  RETURN {
    final_mastery: m_current,
    trajectory: trajectory,
    dropout: dropout,
    steps_completed: len(trajectory),
    steps_total: |P|
  }
```

---

## 4. Path Generation (Phase 2) — Retrieve-Then-Order

**[Defense: #5] Hallucination-proof design**: LLM never generates question
IDs. It only orders questions that were deterministically retrieved.

### 4.1 Step A: Deterministic Retrieval

```
FUNCTION retrieve_candidate_pool(m_n, student_n):
  """
  No LLM involved. Pure rule-based filtering.
  OUTPUT: filtered question pool (typically 30-80 questions)
  """
  pool = []

  # 1. Weak KC questions: mastery < 0.6
  weak_kcs = {c: m_n[c] for c in K if m_n[c] < 0.6}
  FOR c in weak_kcs:
    # ZPD filter: difficulty ∈ [mastery - 0.1, mastery + 0.3]
    zpd_low = max(0.0, m_n[c] - 0.1)
    zpd_high = min(1.0, m_n[c] + 0.3)
    questions = [q for q in Q if c in skill_tags(q)
                               and zpd_low <= difficulty(q) <= zpd_high]
    pool.extend(questions[:10])  # cap per KC to avoid imbalance

  # 2. Stale KC questions: last practiced > 14 days ago
  FOR c in K:
    days = days_since_last_practice(student_n, c)
    IF days > 14 and m_n[c] < 0.8:  # don't review fully mastered
      review_qs = [q for q in Q if c in skill_tags(q)
                                 and difficulty(q) <= m_n[c] + 0.1]
      pool.extend(review_qs[:3])  # fewer review questions

  # 3. Deduplicate
  pool = list(set(pool))

  # 4. Sanity check: need at least L questions
  IF len(pool) < L:
    # Expand ZPD range
    pool = expand_zpd_range(m_n, target_size=L*2)

  RETURN pool
```

### 4.2 Step B: LLM-Based Ordering

```
FUNCTION generate_candidate_paths(m_n, pool, K=5, L=8):
  """
  LLM receives the pool and produces K different orderings of L questions.
  LLM CANNOT hallucinate question_ids — they all come from pool.
  """
  pool_description = format_pool(pool)
  # Format: "q_1234 | KC: fractions, ratios | difficulty: 0.45"

  paths = []
  FOR i = 1 to K:
    prompt = PATH_ORDERING_PROMPT.format(
      mastery_state=format_mastery(m_n),
      question_pool=pool_description,
      path_length=L,
      previous_paths=format_previous(paths) if i > 1 else "None",
      strategy_hint=STRATEGY_HINTS[i]  # different strategy per path
    )
    response = LLM.generate(prompt, temperature=0.6)
    path_i = parse_json_array(response)

    # Validate: all IDs must be in pool
    path_i = [q for q in path_i if q in pool_ids]
    path_i = deduplicate(path_i)[:L]

    # Fill if short (should rarely happen with retrieve-then-order)
    WHILE len(path_i) < L:
      filler = random.choice([q for q in pool if q not in path_i])
      path_i.append(filler)

    paths.append(path_i)

    # LOG for reproducibility [Defense: #6]
    log_llm_call(prompt, response, model_version, temperature)

  RETURN paths
```

### 4.3 Strategy Hints for Diversity

```
STRATEGY_HINTS = {
  1: "Prioritize the weakest KCs first, gradually increasing difficulty.",
  2: "Start with a confidence-building easy question, then tackle weak areas.",
  3: "Interleave weak KCs with stale review KCs throughout the path.",
  4: "Focus on the single weakest KC with intensive practice (3-4 questions).",
  5: "Spread across as many different weak KCs as possible (breadth-first)."
}
```

### 4.4 Path Ordering Prompt Template

```
PATH_ORDERING_PROMPT = """
You are an expert educational path planner. Select and order exactly
{path_length} questions from the pool below to create an optimal
learning path for this student.

## Student Knowledge State
{mastery_state}

## Available Question Pool
{question_pool}

## Strategy
{strategy_hint}

## Constraints
1. Select exactly {path_length} questions from the pool above.
2. Do NOT invent question IDs — use ONLY IDs from the pool.
3. No more than 2 consecutive questions on the same KC (interleaving).
4. Difficulty should have a coherent progression (not random jumps).

## Previously Generated Paths (ensure diversity)
{previous_paths}

## Output
Return ONLY a JSON array of question_ids:
["q_1234", "q_5678", ...]
"""
```

### 4.5 Hallucination Rate Tracking [Defense: #5]

```
METRICS TO REPORT:
  - pool_hit_rate: % of LLM-selected IDs that exist in pool
    (should be ~100% with retrieve-then-order)
  - constraint_violation_rate: % of paths violating interleaving rule
  - fill_rate: % of path slots filled by fallback filler
    (lower = better LLM compliance)

Report these in paper as evidence of generation reliability.
```

---

## 5. Persona Simulation (Phase 3)

### 5.1 Persona Parameter Definitions

```
θ_grit (Gritty Persona):
  learn_rate:              0.15
  forget_penalty:          0.05
  dropout_base:            λ(k) = 0.02 * k       # k = consecutive wrong
  difficulty_sensitivity:  0.2
  skip_threshold:          None                    # never skips
  p_correct_boost:         +0.05
  time_multiplier:         1.3

θ_frag (Fragile Persona):
  learn_rate:              0.10
  forget_penalty:          0.10
  dropout_base:            λ(k) = 0.15 * k^1.5    # aggressive dropout
  difficulty_sensitivity:  0.5                     # very sensitive to hard Qs
  skip_threshold:          0.7                     # skips difficulty > 0.7
  p_correct_boost:         -0.08
  time_multiplier:         0.7

θ_real (Realistic Persona — per-student):
  All parameters extracted from student n's actual data.
  See Section 5.2.
```

**Theoretical grounding [Defense: #2]**:
- Gritty: Duckworth et al. (2007) "Grit: Perseverance and Passion for
  Long-Term Goals" — high persistence, delayed gratification
- Fragile: Bandura (1997) "Self-Efficacy" — low self-efficacy leads to
  task avoidance and early termination under difficulty
- Realistic: Data-driven, no theoretical assumption

### 5.2 Realistic Persona Extraction

```
FUNCTION extract_realistic_params(student_n, history):
  sessions = split_into_sessions(history, gap_minutes=30)

  # Learn rate
  correct_pairs = find_same_kc_consecutive_correct(history)
  learn_rate = clip(mean([Δm for Δm in correct_pairs]), 0.01, 0.30)

  # Forget penalty
  gaps = find_same_kc_with_time_gap(history, min_days=3)
  forget_penalty = clip(mean([max(0, m_before - m_after)
                              for (m_before, m_after, _) in gaps]), 0.01, 0.20)

  # Dropout curve (empirical)
  dropout_counts = {k: {"end": 0, "cont": 0} for k in range(1, 6)}
  FOR session in sessions:
    consec = 0
    FOR i, interaction in enumerate(session):
      IF interaction.correct == 0:
        consec += 1
        IF consec in dropout_counts:
          IF i == len(session) - 1:
            dropout_counts[consec]["end"] += 1
          ELSE:
            dropout_counts[consec]["cont"] += 1
      ELSE:
        consec = 0

  dropout_base = lambda k: (dropout_counts.get(k, {"end":1, "cont":1})["end"] /
                             max(1, sum(dropout_counts.get(k, {"end":1, "cont":1}).values())))

  # Difficulty sensitivity
  difficulty_sensitivity = 0.3  # default, could be estimated

  # Skip threshold
  fast_answers = [q for q in history if q.elapsed_time < 5]
  IF len(fast_answers) > 5:
    skip_threshold = percentile([difficulty(q) for q in fast_answers], 25)
  ELSE:
    skip_threshold = None

  # Time multiplier
  pop_mean = global_mean_elapsed_time()
  time_multiplier = mean(student_n.elapsed_times) / pop_mean

  RETURN θ_real(learn_rate, forget_penalty, dropout_base,
                difficulty_sensitivity, skip_threshold,
                p_correct_boost=0.0, time_multiplier=time_multiplier)
```

### 5.3 Dropout Probability Formula

$$P(\text{dropout} \mid k, d) = \min\left(1.0,\ \lambda_\theta(k) \cdot (1 + \alpha_\theta \cdot d)\right)$$

- $k$ = consecutive wrong answers (including skips)
- $d$ = difficulty of current question ∈ [0, 1]
- $\lambda_\theta(k)$ = persona-specific base curve
- $\alpha_\theta$ = persona-specific difficulty sensitivity

### 5.4 Full Simulation Matrix

```
FUNCTION simulate_all_personas(kt_model, m_n, candidate_paths, θ_real_n):
  """
  Returns score tensor S of shape [K, 3, 3]
  Axes: [path_index, persona_index, metric_index]
  Metrics: [delta_mastery, survival_rate, completion_rate]
  """
  personas = [θ_grit, θ_frag, θ_real_n]
  N_sim = 10  # simulations per (path, persona) for variance reduction
  K = len(candidate_paths)
  S = zeros(K, 3, 3)

  FOR i in range(K):
    FOR j in range(3):
      results = []
      FOR sim in range(N_sim):
        # Fix seed for reproducibility: seed = hash(student_id, i, j, sim)
        set_seed(hash_seed(student_n.id, i, j, sim))
        result = simulate_kt(kt_model, m_n, candidate_paths[i], personas[j])
        results.append(result)

      # Aggregate
      S[i][j][0] = mean([delta_mastery(m_n, r.final_mastery) for r in results])
      S[i][j][1] = 1.0 - mean([float(r.dropout) for r in results])
      S[i][j][2] = mean([r.steps_completed / r.steps_total for r in results])

  RETURN S
```

### 5.5 Delta Mastery (focused on targeted KCs)

$$\Delta\text{mastery} = \frac{1}{|W|} \sum_{c \in W} \max\left(0,\ m_{\text{after}}^c - m_{\text{before}}^c\right)$$

$W$ = KCs targeted by the path (weak ∪ stale KCs). Only counts improvement.

---

## 6. Robust Path Selection (Phase 4)

### 6.1 Composite Score

$$\phi_j(P_i) = w_1 \cdot S_{i,j,0} + w_2 \cdot S_{i,j,1} + w_3 \cdot S_{i,j,2}$$

### 6.2 Dynamic Weights

```
FUNCTION compute_weights(m_n, behavioral_features):
  w1, w2, w3 = 0.5, 0.3, 0.2  # mastery, survival, completion

  # Dropout risk adjustment
  IF behavioral_features.recent_dropout_rate > 0.3:
    w2 += 0.15; w1 -= 0.10; w3 -= 0.05

  # Low mastery adjustment
  IF mean(m_n[weak_kcs]) < 0.3:
    w1 += 0.10; w2 -= 0.05; w3 -= 0.05

  # Normalize
  total = w1 + w2 + w3
  RETURN (w1/total, w2/total, w3/total)
```

### 6.3 Minimax Regret Selection

**Step 1**: Best achievable score per persona:
$$\phi_j^* = \max_{i} \phi_j(P_i)$$

**Step 2**: Regret of path $P_i$ under persona $\theta_j$:
$$R(P_i, \theta_j) = \phi_j^* - \phi_j(P_i) \geq 0$$

**Step 3**: Maximum regret of path $P_i$ across all personas:
$$\bar{R}(P_i) = \max_{j} R(P_i, \theta_j)$$

**Step 4**: Select minimax regret path:
$$P^* = \arg\min_{i} \bar{R}(P_i)$$

**Interpretation [Defense: #4]**: Since personas represent possible
session-level states of the SAME student (not different students),
minimax regret finds the path that performs acceptably regardless of
which mental state the student arrives in today. This is analogous to
robust portfolio optimization under uncertain market conditions.

### 6.4 Algorithm

```
FUNCTION select_robust_path(S, w):
  K, J = S.shape[0], 3
  phi = zeros(K, J)
  FOR i in range(K):
    FOR j in range(J):
      phi[i][j] = w[0]*S[i][j][0] + w[1]*S[i][j][1] + w[2]*S[i][j][2]

  phi_star = [max(phi[:, j]) for j in range(J)]

  regret = zeros(K, J)
  FOR i in range(K):
    FOR j in range(J):
      regret[i][j] = phi_star[j] - phi[i][j]

  max_regret = [max(regret[i, :]) for i in range(K)]
  best_idx = argmin(max_regret)

  RETURN best_idx, regret, phi
```

### 6.5 Alternative Strategies (for Ablation 3)

```
average:           argmax_i mean_j(φ_j(P_i))
maximin:           argmax_i min_j(φ_j(P_i))
hurwicz(α=0.3):   argmax_i [0.3·max_j(φ) + 0.7·min_j(φ)]
hurwicz(α=0.7):   argmax_i [0.7·max_j(φ) + 0.3·min_j(φ)]
weighted_minimax:  same as minimax regret but Realistic persona gets 1.5× weight
                   [Defense: #4 — ablation for reviewer who wants Realistic emphasis]
```

---

## 7. Evaluation Protocol

### 7.1 Held-Out Ground Truth Validation [Defense: #1]

**This is the critical defense against circular evaluation.**

```
PROCEDURE ground_truth_validation(student_n, recommended_path):
  """
  Compare SimPath's PREDICTED mastery gain with ACTUAL performance
  on held-out test data (last 10% of student's interactions).
  """
  # SimPath's prediction
  predicted_gain = simulate_kt(AKT, m_n, recommended_path, θ_real_n).delta_mastery

  # Actual performance on held-out data
  test_interactions = student_n.held_out_data  # last 10%, never seen by SimPath
  actual_kc_performance = {}
  FOR interaction in test_interactions:
    FOR kc in interaction.skill_tags:
      IF kc in targeted_kcs(recommended_path):
        actual_kc_performance.setdefault(kc, []).append(interaction.correct)

  actual_mastery = {kc: mean(responses) for kc, responses in actual_kc_performance.items()}
  actual_gain = mean([actual_mastery.get(c, m_n[c]) - m_n[c]
                      for c in targeted_kcs(recommended_path)])

  RETURN predicted_gain, actual_gain

# Aggregate across all test students:
# Report Pearson correlation between predicted_gains and actual_gains
# Report RMSE between predicted and actual
# This shows simulation fidelity independent of the evaluation KT model
```

### 7.2 Offline Evaluation (Main Metrics)

**IMPORTANT**: Evaluation metrics use SAINT (Model B), NOT AKT (Model A).
AKT is used only for recommendation. [Defense: #1]

```
FUNCTION offline_evaluate(saint_model, test_students, method):
  results = {MG: [], SR: [], CE: [], RI: []}

  FOR student_n in test_students:
    m_n = extract_knowledge_state(saint_model, student_n.recent_history)

    # Generate recommendation using the METHOD being evaluated
    recommended_path = method.recommend(student_n, L=8)

    # Evaluate with ALL THREE personas using SAINT (not AKT!)
    persona_scores = []
    FOR θ in [θ_grit, θ_frag, extract_realistic_params(student_n)]:
      sims = [simulate_kt(saint_model, m_n, recommended_path, θ)
              for _ in range(10)]
      mg = mean([delta_mastery(m_n, s.final_mastery) for s in sims])
      sr = 1.0 - mean([float(s.dropout) for s in sims])
      ce = mean([s.steps_completed / s.steps_total for s in sims])

      results[MG].append(mg)
      results[SR].append(sr)
      results[CE].append(ce)
      persona_scores.append(w1*mg + w2*sr + w3*ce)

    # Robustness Index
    ri = 1.0 - (std(persona_scores) / max(max(persona_scores), 1e-8))
    results[RI].append(ri)

  RETURN {k: (mean(v), std(v)) for k, v in results.items()}
```

### 7.3 Metric Definitions

| Metric | Formula | Range | Better |
|--------|---------|-------|--------|
| Mastery Gain (MG) | $(1/\|W\|)\sum_{c \in W} \max(0, m^c_{\text{after}} - m^c_{\text{before}})$ | [0, 1] | Higher |
| Survival Rate (SR) | $1 - \text{mean}(\mathbb{1}[\text{dropout}])$ | [0, 1] | Higher |
| Completion Eff. (CE) | $\text{steps\_completed} / L$ | [0, 1] | Higher |
| Robustness Index (RI) | $1 - \sigma(\phi_1,\phi_2,\phi_3) / \max(\phi_j)$ | [0, 1] | Higher |

### 7.4 Expert Preference Study (replaces user study) [Defense: #3]

```
DESIGN:
  Participants: 30 educators (K-12 teachers + university instructors)
  Recruitment: education conferences, online teacher communities

  Task A — Path Quality Comparison:
    Show 2 recommended paths (SimPath vs best baseline) side by side.
    Blind labels ("Path A" vs "Path B"), randomized order.
    Ask: "Which path is more educationally appropriate? (A / B / Equal)"
    20 student cases per educator.

  Task B — Explanation Quality:
    Show SimPath's simulation-grounded explanation vs generic rule-based
    explanation for the same recommended path.
    Rate on 5-point Likert:
      1. Usefulness: "Does this help you understand the recommendation?"
      2. Trust: "Do you trust this recommendation based on the explanation?"
      3. Pedagogical soundness: "Is the reasoning educationally valid?"

  Analysis:
    Task A: Binomial test (H0: 50/50 preference)
    Task B: Paired Wilcoxon signed-rank test (SimPath vs baseline explanation)
    Report: preference rate, mean Likert scores ± std, p-values
```

---

## 8. Baselines

### 8.1 Six Baselines

```
B1: DKT + Random
  Train DKT. Identify weak KCs. Select L random questions from weak KCs.
  No ordering optimization.

B2: DKT + Rule-based
  Train DKT. Sort KCs by mastery ascending (weakest first).
  For each KC, select questions with difficulty in [mastery, mastery+0.2].
  Fill path greedily, enforce interleaving.

B3: RL-DKT
  Train DKT. Use DQN agent: state=knowledge_state, action=next_question.
  Reward = delta_mastery per step.
  Train for 50K episodes, ε-greedy (ε=0.1) at test.

B4: LPReKL (KT + LLM generate-and-retrieve)
  KT for state estimation. LLM generates reference exercises.
  Retrieve top-N similar from bank. KT computes knowledge promotion score.
  Iterative refinement.
  Implementation: follow published description + contact original authors.

B5: GenMentor (Multi-agent LLM, no simulation)
  5 LLM agents collaborating (no debate, no simulation stress-test).
  Agents: Gap Identifier, Learner Profiler, Dynamic Simulator,
          Path Scheduler, Content Creator.
  Official code: https://github.com/GeminiLight/gen-mentor

B6: SimPath-NoRobust (Ablation baseline)
  Identical to SimPath but path selection = argmax(average score)
  instead of minimax regret. Tests robustness contribution.
```

### 8.2 Baseline Fairness [Defense: #13]

```
PROTOCOL:
  1. For baselines with public code (GenMentor): use official implementation
  2. For baselines without code (LPReKL): reimplement following paper,
     email original authors for clarification on ambiguous details
  3. ALL baseline reimplementations will be released in our GitHub repo
  4. Use the same LLM (GPT-5.4) for all LLM-based baselines to ensure
     fair comparison (LLM quality is not the variable being tested)
  5. Report: "We contacted the authors of [X] to verify our reimplementation.
     [Response / No response received.]"
```

---

## 9. Ablation Studies

### 9.1 Ablation 1: Number of Personas

```
Variants:
  (a) Realistic only (1 persona)
  (b) Gritty + Fragile (2 personas, no data-driven)
  (c) All 3 (default)
  (d) All 3 + Sprint + Meticulous (5 personas)

Sprint persona (cramming):
  learn_rate=0.20, forget_penalty=0.15, dropout_base=0.05*k
  skip_threshold=0.5, time_multiplier=0.5, p_correct_boost=-0.05

Meticulous persona:
  learn_rate=0.12, forget_penalty=0.03, dropout_base=0.01*k
  skip_threshold=None, time_multiplier=2.0, p_correct_boost=+0.10

Expected finding: 3 personas ≈ 5 personas >> 1 persona.
If confirmed, this justifies the 3-persona default. [Defense: #2]
```

### 9.2 Ablation 2: Number of Candidate Paths (K)

```
K ∈ {3, 5, 8, 10, 15}
Report: MG, SR, RI vs K, with LLM token cost per K
Expected: diminishing returns after K=5-8, cost scales linearly
```

### 9.3 Ablation 3: Selection Strategy

```
Strategies: {minimax_regret, average, maximin, hurwicz_0.3, hurwicz_0.7,
             weighted_minimax_regret}

weighted_minimax_regret: Realistic persona gets weight 1.5,
                         Gritty and Fragile get weight 1.0
[Defense: #4 — tests whether extra weight on Realistic helps]

Expected: minimax_regret best on RI, average best on MG, maximin too
conservative. Report full table to let readers judge.
```

### 9.4 Ablation 4: LLM Backbone

```
Main experiments: GPT-5.4 AND Claude Sonnet 4.6 (both reported)
Ablation: compare the two + discuss differences

Report:
  - Performance metrics (MG, SR, CE, RI) per model
  - Pool hit rate per model (should be ~100% for both)
  - Token cost per model
  - Qualitative: explanation style differences

Note for paper: "Open-source model experiments (Llama, Qwen) are
planned as future work to further validate reproducibility."
[Defense: #6]
```

### 9.5 Ablation 5: Simulation Count (N_sim)

```
N_sim ∈ {1, 5, 10, 20, 50}
Report: metric variance vs N_sim
Expected: variance stabilizes around N_sim=10
```

### 9.6 Ablation 6: KT Model Quality [Defense: #11]

```
Compare SimPath performance when using:
  (a) AKT (high accuracy, ~0.78 AUC) for recommendation
  (b) DKT (lower accuracy, ~0.73 AUC) for recommendation

Both evaluated with SAINT.
Expected: SimPath outperforms baselines with BOTH KT models,
though margins may shrink with lower-quality KT.
This shows robustness to KT model quality.
```

---

## 10. Statistical Testing

```
FOR each (metric, method_pair):
  differences = [simpath_score[i] - baseline_score[i] for i in students]

  # Normality check
  _, p_shapiro = shapiro(differences)

  IF p_shapiro > 0.05:
    stat, p_value = paired_ttest(simpath_scores, baseline_scores)
    test_name = "paired t-test"
  ELSE:
    stat, p_value = wilcoxon(differences)
    test_name = "Wilcoxon signed-rank"

  # Effect size
  cohens_d = mean(differences) / std(differences)

  # Bonferroni correction
  # 6 baselines × 4 metrics = 24 comparisons
  adjusted_alpha = 0.05 / 24  # ≈ 0.002
  significant = p_value < adjusted_alpha

  report(metric, method_pair, test_name, p_value, cohens_d, significant)

FOR ablation:
  Use Kruskal-Wallis + post-hoc Dunn's test with Bonferroni correction
```

---

## 11. Compute Budget

```
Per student (K=5, L=8, N_sim=10, 3 personas):
  Retrieval (Step A):     CPU only, < 0.1s
  LLM path ordering:     5 calls × ~500 tokens = 2,500 tokens
  KT simulation:         5 × 3 × 10 = 150 runs, CPU only, < 1s total
  LLM explanation:       1 call × ~400 tokens = 400 tokens
  Total LLM:             ~2,900 tokens/student
  Total latency:         ~5-8 seconds (with parallel LLM calls)
                         [Defense: #10 — fast enough for session-start batch]

For 1,000 test students:
  Total LLM tokens:      ~2.9M tokens
  GPT-5.4 estimated:     ~$3-5 (price depends on actual pricing)
  Claude Sonnet 4.6:     ~$3-5

Full experiment (main + all ablations, both datasets):
  Estimated: ~50M tokens total, ~$50-100
  GPU time: ~8 hours for KT training (AKT + SAINT × 2 datasets)
  Total wall time: 2-3 days
```

---

## 12. Reproducibility Checklist

```
□ Random seeds: {42, 123, 456, 789, 1024} for 5-run average
□ KT models: save checkpoints, report validation AUC
□ LLM calls: log ALL prompts, responses, model version, params
□ Model versions: exact strings (e.g., "gpt-5.4-2026-03-05",
    "claude-sonnet-4-6")
□ Persona params: θ_grit and θ_frag fixed (reported in paper),
    θ_real varies per student (extraction code provided)
□ Simulation seeds: deterministic via hash(student_id, path_idx,
    persona_idx, sim_idx)
□ Statistical tests: report exact p-values, effect sizes, correction
□ Code: GitHub repo with all code, configs, prompts, sample data
□ Data: EdNet (CC BY-NC 4.0), ASSISTments (public)
□ Persona validation clustering: code + results included
□ Baseline reimplementations: code included, author contact logged
```

---

## 13. Paper Structure & Page Budget (Springer LNAI, 16 pages)

```
Section 1: Introduction (2 pages)
  - Problem: one-size-fits-all recommendations ignore student variability
  - Key insight: personas = session states, not fixed types [Defense: #4]
  - Contribution: simulation-in-the-loop + minimax regret
  - Differentiation from Agent4Edu in paragraph 3 [Defense: #7]

Section 2: Related Work (2 pages)
  - 2.1 Learning Path Recommendation (LPReKL, GenMentor, RL-DKT)
  - 2.2 Student Simulation (Agent4Edu, EduAgent, SimClass)
  - 2.3 Robust Decision-Making (minimax regret, portfolio theory)
  - Table comparing Agent4Edu vs SimPath [Defense: #7]

Section 3: SimPath Framework (4 pages)
  - 3.1 Overview (architecture diagram)
  - 3.2 Retrieve-Then-Order Path Generation [Defense: #5]
  - 3.3 Persona Simulator Design (with theoretical grounding) [Defense: #2]
  - 3.4 Minimax Regret Selection (formulas + worked example)
  - 3.5 Explanation Generation

Section 4: Experiments (4 pages)
  - 4.1 Datasets & Setup
  - 4.2 Persona Validation via Clustering (Figure) [Defense: #2]
  - 4.3 Main Results (2 datasets × 6 baselines × 4 metrics)
  - 4.4 Ground Truth Correlation [Defense: #1]
  - 4.5 Ablation Studies (selection strategy, # personas, K, LLM, KT quality)
  - 4.6 Expert Preference Study results

Section 5: Analysis & Discussion (2.5 pages)
  - 5.1 Case Study: 2-3 student simulation trajectories visualized
  - 5.2 When does robustness matter most? (analysis by student type)
  - 5.3 Domain-specific insights (EdNet vs ASSISTments) [Defense: #9]
  - 5.4 Limitations:
      - Offline evaluation, no online A/B test [Defense: #3]
      - Single-session scope [Defense: #12]
      - KT model dependency (shared limitation) [Defense: #11]
      - Compute latency for real-time use [Defense: #10]

Section 6: Conclusion (0.5 pages)

Appendix (if space / supplementary):
  - Full prompt templates
  - Persona parameter sensitivity analysis
  - Per-KC breakdown of results
  - All LLM call logs (sample)
```

---

## 14. RecSys Reframing (if submitting to RecSys 2026)

```
Title adjustment:
  "SimPath: Robust Learning Resource Recommendation via
   Multi-Persona Simulation Testing"

Framing changes:
  - Emphasize "multi-stakeholder recommendation" angle
  - Frame personas as "user segments with uncertain membership"
  - Minimax regret → "robust recommendation under user uncertainty"
  - Connect to RecSys 2025 Best Paper on "unwanted recommendations"
    (Conformal Risk Control)
  - Emphasize the GENERAL paradigm: simulate-then-select is applicable
    to any recommendation domain where user behavior varies

Additional baseline for RecSys:
  - Standard collaborative filtering + popularity as extra baseline
  - Emphasize that educational domain has MEASURABLE ground truth
    (test scores) unlike movie/music recommendations
```